% VER-2026-001: Cryptographic Compliance Verification
% SendCUIEmail - Verification against REQ-2026-001 Requirements
\documentclass[11pt,letterpaper]{article}

% Page layout
\usepackage[margin=1in]{geometry}
\usepackage{parskip}

% Tables
\usepackage{tabularx}
\usepackage{booktabs}
\usepackage{longtable}
\usepackage{array}

% Code listings
\usepackage{listings}
\usepackage{xcolor}

% Math symbols (for checkmark)
\usepackage{amssymb}

% Hyperlinks
\usepackage{hyperref}
\hypersetup{
    colorlinks=true,
    linkcolor=blue,
    urlcolor=blue,
    citecolor=blue
}

% Headers/footers
\usepackage{fancyhdr}
\pagestyle{fancy}
\fancyhf{}
\rhead{VER-2026-001}
\lhead{SendCUIEmail Cryptographic Compliance}
\rfoot{Page \thepage}

% Code listing style
\definecolor{codebg}{rgb}{0.95,0.95,0.95}
\definecolor{codegreen}{rgb}{0,0.6,0}
\definecolor{codegray}{rgb}{0.5,0.5,0.5}
\definecolor{codeblue}{rgb}{0,0,0.8}

\lstset{
    backgroundcolor=\color{codebg},
    basicstyle=\ttfamily\small,
    breaklines=true,
    captionpos=b,
    commentstyle=\color{codegreen},
    keywordstyle=\color{codeblue},
    stringstyle=\color{codegray},
    numbers=left,
    numberstyle=\tiny\color{codegray},
    numbersep=5pt,
    frame=single,
    framesep=5pt,
    xleftmargin=15pt,
    xrightmargin=5pt
}

% Document metadata
\title{
    \textbf{VER-2026-001} \\
    \Large Cryptographic Compliance Verification \\
    \large SendCUIEmail v0.17.3
}
\author{Verification Document}
\date{January 22, 2026}

\begin{document}

\maketitle

\section*{Document Information}

\begin{tabular}{ll}
\textbf{Document ID:} & VER-2026-001 \\
\textbf{Subject:} & Cryptographic Implementation Compliance Verification \\
\textbf{Requirements:} & REQ-2026-001 v1.1 \\
\textbf{Software Version:} & v0.17.3 \\
\textbf{Git Commit:} & \texttt{2048d2e4508e83ca77b264949775835a3d3aba96} \\
\textbf{Commit Date:} & 2026-01-22 10:12:44 -0600 \\
\textbf{Source File:} & \texttt{Encrypt.ps1} \\
\end{tabular}

\section*{Executive Summary}

This document verifies SendCUIEmail's compliance with REQ-2026-001 Cryptographic Compliance Requirements. Each requirement is traced by its REQ ID, with code evidence demonstrating compliance.

\tableofcontents
\newpage

% ============================================================================
\section{Encryption Algorithm Requirements}
% ============================================================================

\subsection{REQ-1.1: AES Algorithm}

\textbf{Requirement:} The system SHALL use AES (Advanced Encryption Standard) as defined in FIPS 197 for symmetric encryption.

\textbf{Verification Method:} Inspection

\begin{samepage}
\textbf{Evidence (Encrypt.ps1, line 347):}
\begin{lstlisting}[language=bash]
$aes = [System.Security.Cryptography.Aes]::Create()
\end{lstlisting}
\textbf{Result:} PASS $\checkmark$ --- Uses .NET Aes class which implements FIPS 197.
\end{samepage}

\subsection{REQ-1.2: 256-bit Key Length}

\textbf{Requirement:} The system SHALL use a key length of 256 bits (AES-256) for all encryption operations.

\textbf{Verification Method:} Inspection

\begin{samepage}
\textbf{Evidence (Encrypt.ps1, lines 61, 344):}
\begin{lstlisting}[language=bash]
$KEY_SIZE = 32        # bytes (256 bits for AES-256)
...
$key = $keyDeriver.GetBytes($KEY_SIZE)
\end{lstlisting}
\textbf{Result:} PASS $\checkmark$ --- Key is 256 bits (32 bytes).
\end{samepage}

\subsection{REQ-1.3: CBC Mode}

\textbf{Requirement:} The system SHALL use Cipher Block Chaining (CBC) mode or an authenticated encryption mode (GCM, CCM).

\textbf{Verification Method:} Inspection

\begin{samepage}
\textbf{Evidence (Encrypt.ps1, line 348):}
\begin{lstlisting}[language=bash]
$aes.Mode = [System.Security.Cryptography.CipherMode]::CBC
\end{lstlisting}
\textbf{Result:} PASS $\checkmark$ --- Uses CBC mode.
\end{samepage}

\subsection{REQ-1.4: Unique Random IV}

\textbf{Requirement:} The system SHALL generate a unique, random Initialization Vector (IV) for each encryption operation.

\textbf{Verification Method:} Test

\begin{samepage}
\textbf{Evidence (Encrypt.ps1, lines 331-335):}
\begin{lstlisting}[language=bash]
$rng = [System.Security.Cryptography.RandomNumberGenerator]::Create()
$salt = New-Object byte[] $SALT_SIZE
$iv = New-Object byte[] $IV_SIZE
$rng.GetBytes($salt)
$rng.GetBytes($iv)
\end{lstlisting}
\textbf{Test:} Encrypted same file twice, verified IVs differ in output.\\
\textbf{Result:} PASS $\checkmark$ --- New random IV generated for each encryption.
\end{samepage}

\subsection{REQ-1.5: IV Size}

\textbf{Requirement:} The IV SHALL be at least 128 bits (16 bytes) for AES-CBC mode.

\textbf{Verification Method:} Inspection

\begin{samepage}
\textbf{Evidence (Encrypt.ps1, line 62):}
\begin{lstlisting}[language=bash]
$IV_SIZE = 16         # bytes (128 bits for AES block size)
\end{lstlisting}
\textbf{Result:} PASS $\checkmark$ --- IV is exactly 128 bits.
\end{samepage}

\subsection{REQ-1.6: IV Storage}

\textbf{Requirement:} The IV SHALL be stored with or transmitted alongside the ciphertext.

\textbf{Verification Method:} Inspection

\begin{samepage}
\textbf{Evidence (Encrypt.ps1, lines 358-361):}
\begin{lstlisting}[language=bash]
# Combine: Salt + IV + Ciphertext
$outputBytes = New-Object byte[] ($SALT_SIZE + $IV_SIZE + $cipherBytes.Length)
[Array]::Copy($salt, 0, $outputBytes, 0, $SALT_SIZE)
[Array]::Copy($iv, 0, $outputBytes, $SALT_SIZE, $IV_SIZE)
\end{lstlisting}
\textbf{Result:} PASS $\checkmark$ --- IV stored at bytes 16-31 of output file.
\end{samepage}

\newpage
% ============================================================================
\section{Key Derivation Requirements}
% ============================================================================

\subsection{REQ-2.1: PBKDF2}

\textbf{Requirement:} The system SHALL use PBKDF2 (Password-Based Key Derivation Function 2) as defined in NIST SP 800-132 for deriving encryption keys from passwords.

\textbf{Verification Method:} Inspection

\begin{samepage}
\textbf{Evidence (Encrypt.ps1, lines 338-343):}
\begin{lstlisting}[language=bash]
$keyDeriver = New-Object System.Security.Cryptography.Rfc2898DeriveBytes(
    $Password,
    $salt,
    $ITERATIONS,
    [System.Security.Cryptography.HashAlgorithmName]::SHA256
)
\end{lstlisting}
\textbf{Result:} PASS $\checkmark$ --- Uses Rfc2898DeriveBytes (PBKDF2 implementation).
\end{samepage}

\subsection{REQ-2.2: HMAC-SHA256 PRF}

\textbf{Requirement:} The system SHALL use HMAC-SHA256 as the pseudorandom function (PRF) for PBKDF2.

\textbf{Verification Method:} Inspection

\begin{samepage}
\textbf{Evidence (Encrypt.ps1, line 342):}
\begin{lstlisting}[language=bash]
[System.Security.Cryptography.HashAlgorithmName]::SHA256
\end{lstlisting}
\textbf{Result:} PASS $\checkmark$ --- Explicitly specifies SHA256 for HMAC.
\end{samepage}

\subsection{REQ-2.3: Iteration Count}

\textbf{Requirement:} The system SHALL use a minimum of 100,000 iterations for PBKDF2 per NIST SP 800-132 recommendations.

\textbf{Verification Method:} Inspection

\begin{samepage}
\textbf{Evidence (Encrypt.ps1, line 59):}
\begin{lstlisting}[language=bash]
$ITERATIONS = 100000  # PBKDF2 iterations
\end{lstlisting}
\textbf{Result:} PASS $\checkmark$ --- Uses exactly 100,000 iterations.
\end{samepage}

\subsection{REQ-2.4: Unique Random Salt}

\textbf{Requirement:} The system SHALL generate a unique, random salt for each encryption operation.

\textbf{Verification Method:} Test

\begin{samepage}
\textbf{Evidence (Encrypt.ps1, lines 331-334):}
\begin{lstlisting}[language=bash]
$rng = [System.Security.Cryptography.RandomNumberGenerator]::Create()
$salt = New-Object byte[] $SALT_SIZE
...
$rng.GetBytes($salt)
\end{lstlisting}
\textbf{Test:} Encrypted same file twice with same password, verified salts differ.\\
\textbf{Result:} PASS $\checkmark$ --- New random salt generated for each encryption.
\end{samepage}

\subsection{REQ-2.5: Salt Size}

\textbf{Requirement:} The salt SHALL be at least 128 bits (16 bytes) as recommended by NIST SP 800-132.

\textbf{Verification Method:} Inspection

\begin{samepage}
\textbf{Evidence (Encrypt.ps1, line 60):}
\begin{lstlisting}[language=bash]
$SALT_SIZE = 16       # bytes (128 bits)
\end{lstlisting}
\textbf{Result:} PASS $\checkmark$ --- Salt is exactly 128 bits.
\end{samepage}

\subsection{REQ-2.6: Salt Storage}

\textbf{Requirement:} The salt SHALL be stored with or transmitted alongside the ciphertext.

\textbf{Verification Method:} Inspection

\begin{samepage}
\textbf{Evidence (Encrypt.ps1, lines 358-359):}
\begin{lstlisting}[language=bash]
$outputBytes = New-Object byte[] ($SALT_SIZE + $IV_SIZE + $cipherBytes.Length)
[Array]::Copy($salt, 0, $outputBytes, 0, $SALT_SIZE)
\end{lstlisting}
\textbf{Result:} PASS $\checkmark$ --- Salt stored at bytes 0-15 of output file.
\end{samepage}

\subsection{REQ-2.7: Iteration Count Review (Recommended)}

\textbf{Requirement:} The iteration count SHOULD be increased over time as computing power increases.

\textbf{Verification Method:} Analysis

\textbf{Analysis:} Current iteration count (100,000) was set in 2026. NIST recommends reviewing iteration counts periodically. The constant \texttt{\$ITERATIONS} is centrally defined for easy updates.

\textbf{Result:} PASS $\checkmark$ --- Iteration count is configurable and documented.

\newpage
% ============================================================================
\section{Random Number Generation Requirements}
% ============================================================================

\subsection{REQ-3.1: CSPRNG}

\textbf{Requirement:} The system SHALL use a cryptographically secure random number generator (CSPRNG) for generating salts, IVs, and other cryptographic material.

\textbf{Verification Method:} Inspection

\begin{samepage}
\textbf{Evidence (Encrypt.ps1, line 331):}
\begin{lstlisting}[language=bash]
$rng = [System.Security.Cryptography.RandomNumberGenerator]::Create()
\end{lstlisting}
\textbf{Result:} PASS $\checkmark$ --- Uses .NET RandomNumberGenerator (CSPRNG).
\end{samepage}

\subsection{REQ-3.2: Windows System RNG}

\textbf{Requirement:} On Windows, the system SHALL use the system RNG provided by .NET (RNGCryptoServiceProvider or RandomNumberGenerator).

\textbf{Verification Method:} Inspection

\textbf{Evidence:} Same as REQ-3.1. RandomNumberGenerator.Create() returns the platform-appropriate CSPRNG (Windows CNG on Windows).

\textbf{Result:} PASS $\checkmark$

\subsection{REQ-3.3: NIST SP 800-90A Compliance}

\textbf{Requirement:} The CSPRNG SHALL meet the requirements of NIST SP 800-90A or equivalent.

\textbf{Verification Method:} Analysis

\textbf{Analysis:} Windows CNG (bcryptprimitives.dll) is FIPS 140-2 validated (CMVP \#4515) and implements NIST SP 800-90A compliant random number generation.

\textbf{Result:} PASS $\checkmark$

\newpage
% ============================================================================
\section{Password Requirements}
% ============================================================================

\subsection{REQ-4.1: Minimum Password Length}

\textbf{Requirement:} The system SHALL enforce a minimum password length of 8 characters.

\textbf{Verification Method:} Test

\begin{samepage}
\textbf{Evidence (Encrypt.ps1, lines 50, 269-272):}
\begin{lstlisting}[language=bash]
$MIN_PASSWORD_LENGTH = 8
...
if ($plain.Length -lt $MIN_PASSWORD_LENGTH) {
    Write-Host "Password must be at least $MIN_PASSWORD_LENGTH characters"
    ...
\end{lstlisting}
\textbf{Test:} Attempted encryption with 7-character password; rejected.\\
\textbf{Result:} PASS $\checkmark$
\end{samepage}

\subsection{REQ-4.2: Password Confirmation}

\textbf{Requirement:} The system SHALL require password confirmation during encryption to prevent typos.

\textbf{Verification Method:} Test

\begin{samepage}
\textbf{Evidence (Encrypt.ps1, lines 280-293):}
\begin{lstlisting}[language=bash]
$confirm = Read-Host "Confirm password" -AsSecureString
...
if ($plain -ne $confirmPlain) {
    Write-Host "Passwords do not match!" -ForegroundColor Red
\end{lstlisting}
\textbf{Test:} Entered mismatched passwords; rejected with error message.\\
\textbf{Result:} PASS $\checkmark$
\end{samepage}

\subsection{REQ-4.3: Auto Passphrase Generation (Recommended)}

\textbf{Requirement:} The system SHOULD offer automatic passphrase generation using cryptographic RNG.

\textbf{Verification Method:} Test

\begin{samepage}
\textbf{Evidence (Encrypt.ps1, lines 230-244, 258-262):}
\begin{lstlisting}[language=bash]
function New-Passphrase {
    param([int]$WordCount = 4)
    $rng = [System.Security.Cryptography.RandomNumberGenerator]::Create()
    ...
\end{lstlisting}
\textbf{Test:} Pressed Enter at password prompt; received random passphrase.\\
\textbf{Result:} PASS $\checkmark$
\end{samepage}

\subsection{REQ-4.4: No Password Storage}

\textbf{Requirement:} Passwords SHALL NOT be stored or logged by the system.

\textbf{Verification Method:} Inspection

\textbf{Evidence:} Code review confirms password is only held in memory during encryption, then discarded. No file writes or logging of password values.

\textbf{Result:} PASS $\checkmark$

\subsection{REQ-4.5: Masked Password Input}

\textbf{Requirement:} Password input SHALL be masked (not displayed on screen) during entry.

\textbf{Verification Method:} Test

\begin{samepage}
\textbf{Evidence (Encrypt.ps1, line 267):}
\begin{lstlisting}[language=bash]
$password = Read-Host "Enter encryption password" -AsSecureString
\end{lstlisting}
\textbf{Test:} Typed password; characters not displayed (asterisks not shown either per SecureString behavior).\\
\textbf{Result:} PASS $\checkmark$
\end{samepage}

\newpage
% ============================================================================
\section{File Format Requirements}
% ============================================================================

\subsection{REQ-5.1: File Contents}

\textbf{Requirement:} Encrypted files SHALL contain: salt, IV, and ciphertext in a defined format.

\textbf{Verification Method:} Inspection

\textbf{Evidence:} See REQ-5.2 for format specification.

\textbf{Result:} PASS $\checkmark$

\subsection{REQ-5.2: File Format}

\textbf{Requirement:} The file format SHALL be: [16-byte salt][16-byte IV][ciphertext].

\textbf{Verification Method:} Test

\begin{samepage}
\textbf{Evidence (Encrypt.ps1, lines 358-361):}
\begin{lstlisting}[language=bash]
$outputBytes = New-Object byte[] ($SALT_SIZE + $IV_SIZE + $cipherBytes.Length)
[Array]::Copy($salt, 0, $outputBytes, 0, $SALT_SIZE)
[Array]::Copy($iv, 0, $outputBytes, $SALT_SIZE, $IV_SIZE)
[Array]::Copy($cipherBytes, 0, $outputBytes, $SALT_SIZE + $IV_SIZE, ...)
\end{lstlisting}
\textbf{Test:} Hex dump of .Locked file shows 16-byte salt, 16-byte IV, then ciphertext.\\
\textbf{Result:} PASS $\checkmark$
\end{samepage}

\subsection{REQ-5.3: File Extension}

\textbf{Requirement:} Encrypted files SHALL use the .Locked extension to indicate encrypted status.

\textbf{Verification Method:} Inspection

\begin{samepage}
\textbf{Evidence (Encrypt.ps1, line 364):}
\begin{lstlisting}[language=bash]
$outputPath = "$InputPath.Locked"
\end{lstlisting}
\textbf{Result:} PASS $\checkmark$
\end{samepage}

\subsection{REQ-5.4: Original Extension Preserved}

\textbf{Requirement:} The original file extension SHALL be preserved (e.g., document.pdf becomes document.pdf.Locked).

\textbf{Verification Method:} Test

\textbf{Test:} Encrypted \texttt{test.pdf}; output was \texttt{test.pdf.Locked}.

\textbf{Result:} PASS $\checkmark$

\newpage
% ============================================================================
\section{Platform Requirements}
% ============================================================================

\subsection{REQ-6.1: Platform Crypto Libraries}

\textbf{Requirement:} The system SHALL use platform-provided cryptographic libraries (no custom implementations).

\textbf{Verification Method:} Inspection

\textbf{Evidence:} All cryptographic operations use \texttt{System.Security.Cryptography} namespace classes. No custom algorithm implementations.

\textbf{Result:} PASS $\checkmark$

\subsection{REQ-6.2: .NET Cryptographic Classes}

\textbf{Requirement:} On Windows, the system SHALL use .NET Framework cryptographic classes (System.Security.Cryptography).

\textbf{Verification Method:} Inspection

\textbf{Evidence:} Uses:
\begin{itemize}
    \item \texttt{System.Security.Cryptography.Aes}
    \item \texttt{System.Security.Cryptography.Rfc2898DeriveBytes}
    \item \texttt{System.Security.Cryptography.RandomNumberGenerator}
\end{itemize}

\textbf{Result:} PASS $\checkmark$

\subsection{REQ-6.3: FIPS Mode Warning (Recommended)}

\textbf{Requirement:} The system SHOULD detect and warn if FIPS mode is not enabled on Windows.

\textbf{Verification Method:} Test

\begin{samepage}
\textbf{Evidence (Encrypt.ps1, lines 76-108):}
\begin{lstlisting}[language=bash]
function Test-FIPSMode {
    ...
    $fipsKey = "HKLM:\SYSTEM\CurrentControlSet\Control\Lsa\FIPSAlgorithmPolicy"
    ...
\end{lstlisting}
\textbf{Test:} Ran on system with FIPS mode disabled; warning displayed.\\
\textbf{Result:} PASS $\checkmark$
\end{samepage}

\subsection{REQ-6.4: No Third-Party Software}

\textbf{Requirement:} The system SHALL NOT require installation of third-party software on recipient systems.

\textbf{Verification Method:} Analysis

\textbf{Analysis:} Recipients only need Windows PowerShell (built-in) to decrypt files. The Decrypt\_Instructions.html contains a one-liner that works on any Windows system without installation.

\textbf{Result:} PASS $\checkmark$

\newpage
% ============================================================================
\section{Verification Summary}
% ============================================================================

\small
\begin{longtable}{|l|l|l|c|}
\hline
\textbf{REQ ID} & \textbf{Description} & \textbf{Method} & \textbf{Status} \\
\hline
\endhead

\hline
\endfoot

REQ-1.1 & AES Algorithm & Inspection & $\checkmark$ \\
\hline
REQ-1.2 & 256-bit Key & Inspection & $\checkmark$ \\
\hline
REQ-1.3 & CBC Mode & Inspection & $\checkmark$ \\
\hline
REQ-1.4 & Unique Random IV & Test & $\checkmark$ \\
\hline
REQ-1.5 & IV Size (128-bit) & Inspection & $\checkmark$ \\
\hline
REQ-1.6 & IV Storage & Inspection & $\checkmark$ \\
\hline
REQ-2.1 & PBKDF2 & Inspection & $\checkmark$ \\
\hline
REQ-2.2 & HMAC-SHA256 & Inspection & $\checkmark$ \\
\hline
REQ-2.3 & 100k Iterations & Inspection & $\checkmark$ \\
\hline
REQ-2.4 & Unique Random Salt & Test & $\checkmark$ \\
\hline
REQ-2.5 & Salt Size (128-bit) & Inspection & $\checkmark$ \\
\hline
REQ-2.6 & Salt Storage & Inspection & $\checkmark$ \\
\hline
REQ-2.7 & Iteration Review & Analysis & $\checkmark$ \\
\hline
REQ-3.1 & CSPRNG & Inspection & $\checkmark$ \\
\hline
REQ-3.2 & Windows RNG & Inspection & $\checkmark$ \\
\hline
REQ-3.3 & NIST SP 800-90A & Analysis & $\checkmark$ \\
\hline
REQ-4.1 & Min Password Length & Test & $\checkmark$ \\
\hline
REQ-4.2 & Password Confirmation & Test & $\checkmark$ \\
\hline
REQ-4.3 & Auto Passphrase & Test & $\checkmark$ \\
\hline
REQ-4.4 & No Password Storage & Inspection & $\checkmark$ \\
\hline
REQ-4.5 & Masked Input & Test & $\checkmark$ \\
\hline
REQ-5.1 & File Contents & Inspection & $\checkmark$ \\
\hline
REQ-5.2 & File Format & Test & $\checkmark$ \\
\hline
REQ-5.3 & .Locked Extension & Inspection & $\checkmark$ \\
\hline
REQ-5.4 & Extension Preserved & Test & $\checkmark$ \\
\hline
REQ-6.1 & Platform Libraries & Inspection & $\checkmark$ \\
\hline
REQ-6.2 & .NET Crypto Classes & Inspection & $\checkmark$ \\
\hline
REQ-6.3 & FIPS Mode Warning & Test & $\checkmark$ \\
\hline
REQ-6.4 & No Third-Party SW & Analysis & $\checkmark$ \\
\hline

\end{longtable}
\normalsize

\textbf{Summary:} All 28 requirements from REQ-2026-001 v1.1 verified. 25 mandatory requirements PASS, 3 recommended requirements PASS.

% ============================================================================
\section{Conclusion}
% ============================================================================

SendCUIEmail v0.17.3 (commit \texttt{2048d2e}) demonstrates full compliance with REQ-2026-001 v1.1 Cryptographic Compliance Requirements.

\vspace{24pt}

\noindent\rule{\textwidth}{0.5pt}

\textbf{Document prepared:} January 22, 2026 \\
\textbf{Verified against commit:} \texttt{2048d2e4508e83ca77b264949775835a3d3aba96}

\end{document}
